% Mapa visual y glosario de la tesis: panorama matemático con diagramas,
% glosario de terminología (desde lo básico), palabras con varios significados
% y tabla de símbolos. Sin Lean.
% Compilar dos veces:  lualatex MapaVisual.tex
\documentclass[11pt,a4paper]{article}
\usepackage[a4paper,margin=2.2cm,headheight=14pt]{geometry}
\usepackage{amsmath,mathtools}
\usepackage{fontspec}
\usepackage{unicode-math}
\usepackage[spanish,es-noshorthands,es-tabla]{babel}
\setmainfont{TeX Gyre Pagella}
\setmathfont{TeX Gyre Pagella Math}
\setsansfont{DejaVu Sans}[Scale=MatchLowercase]
\usepackage{xcolor}
\usepackage[most]{tcolorbox}
\usepackage{booktabs,tabularx,array,longtable}
\usepackage{bussproofs}
\usepackage{tikz}
\usetikzlibrary{arrows.meta,positioning,trees,fit,calc,shapes.geometric,backgrounds}
\usepackage{tikz-cd}
\usepackage{enumitem}
\usepackage{fancyhdr}
\usepackage[hidelinks]{hyperref}
\usepackage{microtype}
\AtBeginDocument{\providecommand\square{\mdlgwhtsquare}}

\definecolor{c1}{HTML}{3F6FA8}
\definecolor{c2}{HTML}{5E9A62}
\definecolor{c3}{HTML}{D98A4E}
\definecolor{c4}{HTML}{8C7DB8}
\definecolor{c5}{HTML}{C9A227}
\definecolor{fondonota}{HTML}{F1EEF8}
\definecolor{fondoclave}{HTML}{FFF9E6}

\newtcolorbox{nota}[1][Nota]{enhanced, breakable, colback=fondonota,
  colframe=c4, boxrule=0.5pt, arc=3pt, colbacktitle=c4,
  coltitle=white, fonttitle=\sffamily\bfseries\small, title={#1},
  before skip=8pt, after skip=8pt}
\newtcolorbox{clave}{enhanced, colback=fondoclave, colframe=c5,
  boxrule=0.8pt, arc=3pt, left=8pt, right=8pt, top=6pt, bottom=6pt,
  before skip=10pt, after skip=10pt}
\newtcolorbox{figura}[1]{enhanced, colback=white, colframe=c1!50,
  boxrule=0.5pt, arc=3pt, colbacktitle=c1!12, coltitle=black,
  fonttitle=\sffamily\small\bfseries, title={#1}, halign=center,
  before skip=8pt, after skip=10pt}

\newcommand{\Provnd}{\vdash_{\mathrm{ND}}}
\newcommand{\sem}[1]{[\![#1]\!]}
\newcommand{\bred}{\rightarrow_\beta}
\newcommand{\bmred}{\twoheadrightarrow_\beta}
\newcommand{\concat}{\mathbin{+\!\!+}}
\newcommand{\qs}{\operatorname{\mathsf{qsort}}}
\newcommand{\termino}[1]{\textbf{#1}}

\tikzset{
  caja/.style={draw, rounded corners, align=center, inner sep=4pt, font=\small},
  azul/.style={caja, fill=c1!10, draw=c1},
  verde/.style={caja, fill=c2!12, draw=c2},
  naranja/.style={caja, fill=c3!12, draw=c3},
  morado/.style={caja, fill=c4!12, draw=c4},
  amarillo/.style={caja, fill=c5!15, draw=c5},
  flecha/.style={-{Stealth}, thick, gray!60!black},
  dflecha/.style={{Stealth}-{Stealth}, thick, gray!60!black},
  et/.style={font=\scriptsize, fill=white, inner sep=1pt},
}

\pagestyle{fancy}
\fancyhf{}
\fancyhead[L]{\small\sffamily Mapa visual y glosario de la tesis}
\fancyhead[R]{\small\sffamily\nouppercase{\leftmark}}
\fancyfoot[C]{\small\thepage}
\renewcommand{\headrulewidth}{0.4pt}
\setlist{itemsep=2pt, topsep=4pt}
\setlength{\parskip}{4pt}
\setlength{\parindent}{0pt}
\EnableBpAbbreviations

\title{\textbf{Mapa visual y glosario de la tesis}\\[4pt]
  \large Panorama matemático en diagramas y la terminología\\ que hay que dominar, desde lo básico}
\author{Material de apoyo para la defensa}
\date{}

\begin{document}
\maketitle
\thispagestyle{empty}

\begin{abstract}
\noindent
Este documento es para tener \emph{fluidez}: ver la tesis completa de un
vistazo y poder usar cada término con precisión. La primera parte recorre las
ideas matemáticas con diagramas, desde la sintaxis y la semántica hasta los
dos teoremas principales. La segunda es un glosario alfabético que incluye
también los términos básicos (relación, inducción, base, función…). La
tercera aclara palabras que la tesis usa con más de un significado
(«completo», «corrección», «base», «normal», «\(\vdash\)»…). Cierra una tabla de
símbolos. No usa código.
\end{abstract}

\tableofcontents
\newpage

%=====================================================================
\section{La tesis en una imagen}
%=====================================================================

\begin{figura}{El recorrido completo: de la teoría a los dos estudios de caso}
\begin{tikzpicture}[node distance=6mm and 6mm]
\node[azul, text width=34mm] (c7) {\textbf{Cap.\ 7}\\cálculo \(\lambda\)\\y tipos simples};
\node[azul, text width=34mm, right=of c7] (c8) {\textbf{Cap.\ 8}\\cubo \(\lambda\)\\hasta el CoC};
\node[verde, text width=34mm, right=of c8] (c9) {\textbf{Cap.\ 9}\\deducción natural,\\semántica, Curry--Howard};
\node[azul, text width=34mm, below=12mm of c8] (c10) {\textbf{Cap.\ 10}\\CIC: tipos inductivos\\y eliminadores};
\node[naranja, text width=42mm, below left=14mm and -8mm of c10] (c11) {\textbf{Cap.\ 11}\\\(\mathrm{Taut}(\varphi)\Rightarrow\varnothing\Provnd\varphi\)\\(completitud)};
\node[naranja, text width=42mm, below right=14mm and -8mm of c10] (c12) {\textbf{Cap.\ 12}\\\(\qs(\ell)\sim\ell\ \wedge\ \mathrm{Ordenada}(\qs(\ell))\)\\(Quicksort)};
\node[morado, text width=52mm, below=34mm of c10] (lean) {\textbf{Asistente de pruebas}\\el núcleo comprueba cada prueba;\\auditoría de axiomas};
\draw[flecha] (c7) -- node[et,above]{tipos} (c8);
\draw[flecha] (c8) -- node[et,above]{lógica} (c9);
\draw[flecha] (c8) -- node[et,right]{+ inductivos} (c10);
\draw[flecha] (c9) |- node[et,pos=0.25,right]{qué se formaliza} (c10);
\draw[flecha] (c10) -- (c11);
\draw[flecha] (c10) -- (c12);
\draw[flecha] (c11) -- (lean);
\draw[flecha] (c12) -- (lean);
\end{tikzpicture}
\end{figura}

\subsection*{La tesis en diez frases}
\begin{enumerate}
  \item El \termino{cálculo \(\lambda\)} es un lenguaje mínimo de funciones:
    variables, aplicación y abstracción. Calcular es \termino{reducir}.
  \item Los \termino{tipos} impiden expresiones sin sentido (\(x\,x\)) y
    garantizan que el cálculo termina.
  \item El \termino{cubo \(\lambda\)} añade tres tipos de dependencia; en el
    vértice más fuerte está el \termino{Cálculo de Construcciones} (CoC).
  \item La \termino{correspondencia de Curry--Howard}: proposiciones son tipos
    y demostraciones son programas.
  \item El \termino{CIC} añade al CoC los \termino{tipos inductivos}; con ellos
    vienen gratis la recursión y la inducción (eliminadores).
  \item Un \termino{asistente de pruebas} implementa el CIC: si un término
    tiene el tipo de un enunciado, ese término es una demostración comprobada.
  \item La \termino{deducción natural} define qué es demostrar (\(\vdash\)); las
    \termino{tablas de verdad} definen qué es verdadero (\(\models\)).
  \item \termino{Completitud}: todo lo verdadero en todas las filas es
    demostrable. Se prueba con el \termino{argumento de Kalmár}.
  \item \termino{Quicksort} termina y devuelve una permutación ordenada de su
    entrada; se prueba por inducción sobre la longitud.
  \item La \termino{auditoría de axiomas} certifica que no hay huecos y que solo
    se usan axiomas estándar.
\end{enumerate}

%=====================================================================
\section{Sintaxis y semántica}
%=====================================================================

\begin{figura}{Los dos mundos y los teoremas que los unen}
\begin{tikzpicture}
\node[azul, text width=48mm, minimum height=30mm] (sin) at (0,0) {\textbf{Sintaxis}\\[3pt] símbolos y reglas\\[3pt]\(\Gamma\vdash\varphi\)\\«existe una derivación»\\[3pt]\footnotesize deducción natural, árboles, cajas};
\node[verde, text width=48mm, minimum height=30mm] (sem) at (10.0,0) {\textbf{Semántica}\\[3pt] significado y verdad\\[3pt]\(\Gamma\models\varphi\)\\«toda valuación que hace verdadero \(\Gamma\) hace verdadera \(\varphi\)»\\[3pt]\footnotesize tablas de verdad};
\draw[flecha, very thick] ([yshift=6mm]sin.east) -- node[et,above,align=center]{\textbf{corrección}\\(solidez)\ \(\vdash\Rightarrow\models\)} ([yshift=6mm]sem.west);
\draw[flecha, very thick] ([yshift=-6mm]sem.west) -- node[et,below,align=center]{\textbf{completitud}\\\(\models\Rightarrow\vdash\)} ([yshift=-6mm]sin.east);
\node[amarillo] at (5.2,-2.6) {\textbf{adecuación}: \(\ \vdash\ \Longleftrightarrow\ \models\)};
\end{tikzpicture}
\end{figura}

\begin{itemize}
  \item \termino{Corrección} dice que las reglas no se equivocan: lo
    demostrable es verdadero. Es la dirección «fácil».
  \item \termino{Completitud} dice que las reglas no se quedan cortas: lo
    verdadero es demostrable. Es la dirección difícil y el Capítulo 11.
  \item \termino{Consistencia} es un caso particular de la corrección: como
    \(\bot\) nunca es verdadera, nunca es demostrable.
\end{itemize}

\begin{figura}{Clasificación de las fórmulas por su tabla de verdad}
\begin{tikzpicture}
\draw[thick, rounded corners, fill=gray!5] (0,0) rectangle (12,4.2);
\node[anchor=north west, font=\small\bfseries] at (0.1,4.15) {todas las fórmulas};
\draw[thick, rounded corners, fill=c2!10, draw=c2] (0.4,0.4) rectangle (8.2,3.6);
\node[anchor=north west, font=\small] at (0.5,3.55) {\textbf{satisfacibles}: verdaderas en \emph{alguna} fila};
\draw[thick, rounded corners, fill=c2!25, draw=c2] (0.8,0.7) rectangle (4.2,2.8);
\node[align=center, font=\small] at (2.5,1.75) {\textbf{tautologías}\\verdaderas en\\\emph{todas} las filas\\\(p\vee\neg p\),\ \(p\to p\)};
\node[align=center, font=\small] at (6.2,1.75) {contingentes\\\(p\wedge\neg q\)};
\draw[thick, rounded corners, fill=c3!15, draw=c3] (8.6,0.4) rectangle (11.6,3.6);
\node[align=center, font=\small] at (10.1,2) {\textbf{contradicciones}\\falsas en todas\\las filas\\\(p\wedge\neg p\),\ \(\bot\)};
\end{tikzpicture}
\end{figura}

%=====================================================================
\section{Nivel objeto y nivel meta}
%=====================================================================

\begin{figura}{Se estudia un sistema (objeto) razonando en otro (meta)}
\begin{tikzpicture}
\draw[rounded corners, thick, draw=c4, fill=c4!6] (0,0) rectangle (13,5);
\node[anchor=north west, font=\small\bfseries, text=c4] at (0.1,4.95) {METALENGUAJE: matemática ordinaria (en la tesis, la lógica del asistente)};
\node[font=\small, align=left, anchor=west] at (0.4,4.1) {aquí viven los \emph{teoremas sobre} el cálculo: «si \(\Gamma\vdash\varphi\), entonces \dots», inducción, «por contradicción»};
\draw[rounded corners, thick, draw=c1, fill=c1!8] (0.6,0.5) rectangle (12.4,3.4);
\node[anchor=north west, font=\small\bfseries, text=c1] at (0.7,3.35) {LENGUAJE OBJETO: el cálculo \(\Provnd\)};
\node[font=\small, align=left, anchor=west] at (1,2.45) {fórmulas \(p,\ \neg\varphi,\ \varphi\to\psi\) \quad reglas hip., \(\to\)I, \(\to\)E, \(\neg\)I, \(\neg\)E, \(\bot\)E, RAA};
\node[font=\small, align=left, anchor=west] at (1,1.6) {\(\neg\) es un \emph{símbolo}; \(\bot:=\neg(p_0\to p_0)\) es una \emph{fórmula}};
\node[font=\small, align=left, anchor=west] at (1,0.9) {RAA es una \emph{regla} del cálculo};
\end{tikzpicture}
\end{figura}

\begin{center}
\renewcommand{\arraystretch}{1.25}
{\small
\begin{tabular}{@{}lll@{}}
\toprule
& \textbf{Objeto} & \textbf{Meta}\\
\midrule
negación & \(\neg\varphi\) (una fórmula) & «no \(P\)» (una afirmación)\\
falsedad & \(\bot\) (una fórmula) & «Falso» (una proposición)\\
implicación & \(\varphi\to\psi\) & «si \dots\ entonces \dots» (\(\Rightarrow\))\\
contradicción & regla RAA & razonar por contradicción\\
\bottomrule
\end{tabular}}
\end{center}

%=====================================================================
\section{Objetos definidos por reglas}
%=====================================================================

Casi todo en la tesis se define \termino{inductivamente}: se dan unos
\termino{constructores} y se toma el \emph{menor} conjunto cerrado bajo ellos.
Cada objeto así definido es un \termino{árbol}.

\begin{figura}{El árbol sintáctico de una fórmula y el de un término \(\lambda\)}
\begin{tikzpicture}[level distance=11mm,
  every node/.style={circle, draw=c1, fill=c1!10, minimum size=7mm, inner sep=1pt, font=\small},
  level 1/.style={sibling distance=22mm}, level 2/.style={sibling distance=12mm}]
\node {\(\to\)}
  child {node {\(p\)}}
  child {node {\(\to\)}
    child {node {\(q\)}}
    child {node {\(p\)}}};
\node[draw=none, fill=none, font=\small] at (0,-3.6) {\(p\to(q\to p)\)};
\begin{scope}[xshift=7cm, every node/.style={circle, draw=c2, fill=c2!12, minimum size=7mm, inner sep=1pt, font=\small}]
\node {@}
  child {node {\(\lambda x\)}
    child {node {@}
      child {node {\(x\)}}
      child {node {\(y\)}}}}
  child {node {\(z\)}};
\node[draw=none, fill=none, font=\small] at (0,-3.6) {\((\lambda x.\,x\,y)\,z\)\quad(@ = aplicación)};
\end{scope}
\end{tikzpicture}
\end{figura}

\begin{figura}{Una definición inductiva sirve para dos cosas}
\begin{tikzpicture}
\node[amarillo, text width=40mm] (def) at (0,0) {\textbf{Definición inductiva}\\constructores +\\«nada más»};
\node[azul, text width=50mm] (con) at (-4,-2.3) {\textbf{Construir}\\aplicar reglas finitas veces\\(una derivación, un término)};
\node[verde, text width=50mm] (raz) at (4,-2.3) {\textbf{Razonar sobre todos}\\si una propiedad se conserva\\por cada regla, vale siempre\\(\textbf{inducción})};
\draw[flecha] (def) -- (con);
\draw[flecha] (def) -- (raz);
\end{tikzpicture}
\end{figura}

%=====================================================================
\section{Las formas de inducción}
%=====================================================================

\begin{figura}{Cuatro inducciones, una misma idea}
\begin{tikzpicture}
\node[azul, text width=36mm, minimum height=24mm] (a) at (0,0) {\textbf{Sobre \(\mathbb N\)}\\\(P(0)\)\\\(P(n)\Rightarrow P(n{+}1)\)\\[2pt]\footnotesize dominó};
\node[azul, text width=36mm, minimum height=24mm] (b) at (4.6,0) {\textbf{Estructural}\\un caso por constructor;\\hipótesis para los hijos\\[2pt]\footnotesize fórmulas, términos, listas};
\node[verde, text width=36mm, minimum height=24mm] (c) at (0,-3) {\textbf{Sobre derivaciones}\\un caso por regla;\\hipótesis para las premisas\\[2pt]\footnotesize debilitamiento, corrección};
\node[naranja, text width=36mm, minimum height=24mm] (d) at (4.6,-3) {\textbf{Bien fundada}\\suponer \(P(s)\) para todo\\\(s\) «más pequeño»\\[2pt]\footnotesize Quicksort (longitud)};
\node[morado, text width=52mm] (e) at (10.6,-1.5) {Todas son el \textbf{eliminador} de un tipo inductivo (Cap.\ 10).\\[3pt]La bien fundada vale porque en \((\mathbb N,<)\) no hay cadenas descendentes infinitas.};
\draw[flecha] (b) -- (e); \draw[flecha] (d) -- (e);
\end{tikzpicture}
\end{figura}

\begin{figura}{Por qué la inducción bien fundada es válida}
\begin{tikzpicture}
\foreach \x/\n in {0/6,2/3,4/1,6/0} {\node[azul, minimum width=10mm] (n\n) at (\x,0) {\(\n\)};}
\draw[flecha] (n6) -- node[et,above]{llamada} (n3);
\draw[flecha] (n3) -- (n1);
\draw[flecha] (n1) -- (n0);
\node[font=\small, align=left, anchor=west] at (7.2,0) {toda cadena que baja en \(\mathbb N\)\\termina en finitos pasos};
\node[naranja, font=\small] at (3,-1.3) {en \(\mathbb Z\) no: \(0>-1>-2>\cdots\) nunca para \(\Rightarrow\) \((\mathbb Z,<)\) no es bien fundado};
\end{tikzpicture}
\end{figura}

\begin{figura}{Generalizar antes de inducir}
\begin{tikzpicture}
\node[naranja, text width=58mm] (mal) at (0,0) {\textbf{Mal}: fijar \(\Delta\) y luego inducir.\\En \(\to\)I la premisa tiene contexto \(\Gamma\cup\{\alpha\}\): la hipótesis inductiva habla de \(\Delta\), no de \(\Delta\cup\{\alpha\}\). Se atasca.};
\node[verde, text width=58mm] (bien) at (7.5,0) {\textbf{Bien}: inducir sobre\\\(P(\Gamma,\varphi):=\forall\Delta\supseteq\Gamma,\ \Delta\vdash\varphi\).\\La hipótesis se usa con \(\Delta\cup\{\alpha\}\).};
\draw[flecha] (mal) -- (bien);
\end{tikzpicture}
\end{figura}

%=====================================================================
\section{Relaciones y órdenes}
%=====================================================================

\begin{figura}{Propiedades de una relación, en dibujos}
\begin{tikzpicture}[pt/.style={circle, fill=c1, inner sep=1.8pt}]
\begin{scope}
\node[pt] (a) at (0,0) {}; \draw[flecha] (a) edge[loop above] (a);
\node[font=\small] at (0,-0.8) {reflexiva: \(a\,R\,a\)};
\end{scope}
\begin{scope}[xshift=3.3cm]
\node[pt] (a) at (0,0) {}; \node[pt] (b) at (1.2,0) {};
\draw[flecha] (a) to[bend left] (b); \draw[flecha] (b) to[bend left] (a);
\node[font=\small] at (0.6,-0.8) {simétrica};
\end{scope}
\begin{scope}[xshift=6.8cm]
\node[pt] (a) at (0,0) {}; \node[pt] (b) at (0.7,0.9) {}; \node[pt] (c) at (1.4,0) {};
\draw[flecha] (a) -- (b); \draw[flecha] (b) -- (c); \draw[flecha, dashed, c3] (a) -- (c);
\node[font=\small] at (0.7,-0.8) {transitiva};
\end{scope}
\begin{scope}[xshift=10.3cm]
\node[pt] (a) at (0,0) {}; \node[pt] (b) at (1.2,0) {};
\draw[flecha] (a) -- (b);
\node[font=\small, align=center] at (0.6,-0.9) {antisimétrica: ida y\\vuelta solo si \(a=b\)};
\end{scope}
\end{tikzpicture}
\end{figura}

\begin{figura}{La familia de los órdenes y la de las equivalencias}
\begin{tikzpicture}
\node[azul, text width=40mm] (pre) at (0,0) {\textbf{preorden}\\reflexivo + transitivo};
\node[azul, text width=40mm] (par) at (-3,-1.7) {\textbf{orden parcial}\\+ antisimétrico};
\node[azul, text width=40mm] (tot) at (3,-1.7) {\textbf{preorden total}\\+ total (\(a\leq b\) o \(b\leq a\))};
\node[amarillo, text width=46mm] (lin) at (0,-3.4) {\textbf{orden lineal (total)}\\las cuatro propiedades\\\footnotesize hipótesis de Quicksort};
\node[verde, text width=44mm] (eq) at (8,-0.6) {\textbf{relación de equivalencia}\\reflexiva + simétrica + transitiva\\\footnotesize \(=_\alpha\),\ \(=_\beta\),\ \(\sim\) (permutación)};
\draw[flecha] (pre) -- (par); \draw[flecha] (pre) -- (tot);
\draw[flecha] (par) -- (lin); \draw[flecha] (tot) -- (lin);
\end{tikzpicture}
\end{figura}

\begin{figura}{Clausuras: de un paso a muchos pasos}
\begin{tikzpicture}
\node[azul] (a) at (0,0) {\(\bred\)\\un paso};
\node[azul] (b) at (4.2,0) {\(\bmred\)\\cero o más pasos\\\footnotesize clausura reflexiva y transitiva};
\node[verde] (c) at (9.2,0) {\(=_\beta\)\\pasos en ambos sentidos\\\footnotesize menor equivalencia que contiene \(\bred\)};
\draw[flecha] (a) -- (b); \draw[flecha] (b) -- (c);
\end{tikzpicture}
\end{figura}

%=====================================================================
\section{Cálculo lambda}
%=====================================================================

\begin{figura}{Variables libres y ligadas en \(\lambda x.\,(x\,y)\)}
\begin{tikzpicture}
\node[font=\Large] (t) at (0,0) {\(\lambda\,\underset{\text{ligadora}}{x}\;.\;(\,\underset{\text{ligada}}{x}\;\;\underset{\text{libre}}{y}\,)\)};
\draw[flecha, c2, thick] (-0.95,0.3) to[bend left=50] node[et,above]{enlaza} (0.35,0.3);
\node[font=\small, align=center] at (0,-1.4) {\(\mathrm{FV}(\lambda x.\,x\,y)=\{y\}\). El nombre de \(x\) no importa: \(\lambda x.\,x\,y=_\alpha\lambda u.\,u\,y\).\\Renombrar \(x\) a \(y\) estaría prohibido: \textbf{capturaría} la \(y\) libre.};
\end{tikzpicture}
\end{figura}

\begin{figura}{Un paso \(\beta\): sustituir el argumento en el cuerpo}
\begin{tikzpicture}
\node[azul] (r) at (0,0) {\((\lambda x.\,\underbrace{x\,(x\,y)}_{\text{cuerpo}})\ \underbrace{N}_{\text{argumento}}\)\\\footnotesize redex};
\node[verde] (c) at (6.5,0) {\(N\,(N\,y)\)\\\footnotesize contractum \(=\) cuerpo\([x:=N]\)};
\draw[flecha] (r) -- node[et,above]{\(\bred\)} (c);
\end{tikzpicture}
\end{figura}

\begin{figura}{Confluencia (Church--Rosser) y su ejemplo de la tesis}
\begin{tikzpicture}
\node (m) at (0,0) {\(M\)};
\node (n1) at (-1.8,-1.6) {\(N_1\)};
\node (n2) at (1.8,-1.6) {\(N_2\)};
\node (l) at (0,-3.2) {\(L\)};
\draw[-{Stealth[length=2mm] Stealth[length=2mm]}] (m) -- (n1);
\draw[-{Stealth[length=2mm] Stealth[length=2mm]}] (m) -- (n2);
\draw[-{Stealth[length=2mm] Stealth[length=2mm]}, dashed] (n1) -- (l);
\draw[-{Stealth[length=2mm] Stealth[length=2mm]}, dashed] (n2) -- (l);
\node[font=\small, align=center] at (0,-4) {si se divergen, se vuelven a encontrar};
\begin{scope}[xshift=8cm]
\node (m) at (0,0) {\((\lambda x.(\lambda y.\,y\,x)\,z)\,v\)};
\node (n1) at (-2.2,-1.6) {\((\lambda y.\,y\,v)\,z\)};
\node (n2) at (2.2,-1.6) {\((\lambda x.\,z\,x)\,v\)};
\node (l) at (0,-3.2) {\(z\,v\)};
\draw[flecha] (m) -- node[et,left]{externo} (n1);
\draw[flecha] (m) -- node[et,right]{interno} (n2);
\draw[flecha] (n1) -- (l); \draw[flecha] (n2) -- (l);
\end{scope}
\end{tikzpicture}
\end{figura}

\begin{figura}{Terminación: tres comportamientos}
\begin{tikzpicture}
\node[verde, text width=36mm] (a) at (0,0) {\textbf{fuertemente normalizable}\\toda ruta termina\\\footnotesize todo término tipado};
\node[amarillo, text width=36mm] (b) at (4.6,0) {\textbf{débilmente normalizable}\\alguna ruta termina\\\footnotesize \((\lambda y.\,z)\,\Omega\)};
\node[naranja, text width=36mm] (c) at (9.2,0) {\textbf{no normalizable}\\ninguna ruta termina\\\footnotesize \(\Omega=(\lambda x.xx)(\lambda x.xx)\)\\\footnotesize \(\Omega\bred\Omega\bred\cdots\)};
\end{tikzpicture}
\end{figura}

%=====================================================================
\section{Tipos}
%=====================================================================

\begin{figura}{Anatomía de un juicio de tipado}
\begin{tikzpicture}
\node[font=\Large] at (0,0) {\(\underbrace{x:\sigma,\ y:\tau}_{\text{contexto }\Gamma}\ \ \vdash\ \ \underbrace{\lambda z{:}\sigma.\,x}_{\text{término (sujeto)}}\ :\ \underbrace{\sigma\to\sigma}_{\text{tipo}}\)};
\node[font=\small, text width=140mm, align=center] at (0,-1.5) {«con estas declaraciones, este término tiene este tipo». Bajo Curry--Howard: «con estas hipótesis, esta prueba demuestra esta proposición».};
\end{tikzpicture}
\end{figura}

\begin{figura}{La escalera de clasificación}
\begin{tikzpicture}
\node[verde] (t) at (0,0) {\(3\)\\\footnotesize término};
\node[azul] (ty) at (3,0) {\(\mathsf{nat}\)\\\footnotesize tipo};
\node[morado] (k) at (6,0) {\(*\)\\\footnotesize kind (tipo de los tipos)};
\node[amarillo] (s) at (9.6,0) {\(\square\)\\\footnotesize sort superior};
\draw[flecha] (t) -- node[et,above]{:} (ty);
\draw[flecha] (ty) -- node[et,above]{:} (k);
\draw[flecha] (k) -- node[et,above]{:} (s);
\node[verde] (t2) at (0,-1.6) {\(\lambda\alpha{:}*.\,\alpha\to\alpha\)\\\footnotesize constructor};
\node[morado] (k2) at (4.2,-1.6) {\(*\to*\)\\\footnotesize kind};
\node[amarillo] (s2) at (7.6,-1.6) {\(\square\)};
\draw[flecha] (t2) -- node[et,above]{:} (k2);
\draw[flecha] (k2) -- node[et,above]{:} (s2);
\end{tikzpicture}
\end{figura}

\begin{figura}{El cubo \(\lambda\): tres ejes de dependencia}
\begin{tikzpicture}[scale=1.25, v/.style={azul, font=\small, inner sep=2pt}]
\coordinate (O) at (0,0); \coordinate (X) at (4,0); \coordinate (Y) at (0,3); \coordinate (Z) at (1.6,1.3);
\node[v] (arr) at (O) {\(\lambda\!\to\)};
\node[v] (P) at (X) {\(\lambda P\)};
\node[v] (l2) at (Y) {\(\lambda 2\)};
\node[v] (P2) at ($(X)+(Y)$) {\(\lambda P2\)};
\node[v] (lw) at (Z) {\(\lambda\underline{\omega}\)};
\node[v] (Pw) at ($(X)+(Z)$) {\(\lambda P\underline{\omega}\)};
\node[v] (w) at ($(Y)+(Z)$) {\(\lambda\omega\)};
\node[amarillo, font=\small, inner sep=2pt] (C) at ($(X)+(Y)+(Z)$) {\(\lambda C\) = CoC};
\draw (arr)--(P) (arr)--(l2) (P)--(P2) (l2)--(P2) (l2)--(w) (P2)--(C) (w)--(C) (P)--(Pw) (Pw)--(C);
\draw[dashed] (arr)--(lw) (lw)--(Pw) (lw)--(w);
\node[font=\scriptsize, align=left, anchor=west] at (6.2,2.6) {\(\uparrow\) \textbf{términos que dependen de tipos}\\\quad \((\square,*)\): polimorfismo, \(\lambda 2\)};
\node[font=\scriptsize, align=left, anchor=west] at (6.2,1.5) {\(\nearrow\) \textbf{tipos que dependen de tipos}\\\quad \((\square,\square)\): constructores, \(\lambda\underline{\omega}\)};
\node[font=\scriptsize, align=left, anchor=west] at (6.2,0.4) {\(\rightarrow\) \textbf{tipos que dependen de términos}\\\quad \((*,\square)\): lógica de predicados, \(\lambda P\)};
\end{tikzpicture}
\end{figura}

%=====================================================================
\section{Curry--Howard: pruebas como programas}
%=====================================================================

\begin{figura}{La misma estructura, leída dos veces (Ejemplo 9.24)}
\begin{tikzpicture}
\node at (0,0) {
\AxiomC{\([B\wedge A]^z\)}\UnaryInfC{\(A\)}
\AxiomC{\([B\wedge A]^z\)}\UnaryInfC{\(B\)}
\BinaryInfC{\(A\wedge B\)}
\RightLabel{\scriptsize\(\to\)I}
\UnaryInfC{\((B\wedge A)\to(A\wedge B)\)}
\DisplayProof};
\node at (7.6,0) {
\AxiomC{\([z:B\times A]\)}\UnaryInfC{\(\pi_2 z:A\)}
\AxiomC{\([z:B\times A]\)}\UnaryInfC{\(\pi_1 z:B\)}
\BinaryInfC{\(\langle\pi_2z,\pi_1z\rangle:A\times B\)}
\RightLabel{\scriptsize\(\lambda\)}
\UnaryInfC{\(\lambda z.\langle\pi_2z,\pi_1z\rangle\)}
\DisplayProof};
\node[font=\small\bfseries] at (0,-1.6) {prueba lógica};
\node[font=\small\bfseries] at (7.6,-1.6) {programa (intercambia un par)};
\draw[dflecha] (2.3,-0.2) -- node[et,above]{borrar / añadir términos} (4.9,-0.2);
\end{tikzpicture}
\end{figura}

\begin{center}
\renewcommand{\arraystretch}{1.2}
{\small
\begin{tabular}{@{}ll@{\qquad}ll@{}}
\toprule
\textbf{Lógica} & \textbf{Cómputo} & \textbf{Lógica} & \textbf{Cómputo}\\
\midrule
proposición & tipo & \(A\to B\) & función \(A\to B\)\\
prueba & término (programa) & \(A\wedge B\) & par \(A\times B\)\\
hipótesis & variable & \(A\vee B\) & suma \(A+B\)\\
\(\to\)I (descargar) & \(\lambda\) (ligar) & \(\forall x.\,P(x)\) & \(\Pi x.\,P\,x\)\\
\(\to\)E (modus ponens) & aplicación & \(\bot\) & tipo vacío\\
desvío en una prueba & redex & normalizar & evaluar\\
\bottomrule
\end{tabular}}
\end{center}

%=====================================================================
\section{Tipos de lógica en la tesis}
%=====================================================================

La tesis nombra varias lógicas: clásica, intuicionista, mínima, de Heyting,
proposicional, de predicados, de primer orden, de orden superior y «la lógica
de Lean». No son una lista plana: responden a dos preguntas independientes.

\begin{clave}
\textbf{Pregunta 1: ¿qué cuenta como prueba?} Es decir, qué reglas se
permiten. De menos a más reglas:
\termino{mínima} \(\subset\) \termino{intuicionista} \(\subset\) \termino{clásica}.

\smallskip
\textbf{Pregunta 2: ¿de qué se puede hablar?} Es decir, sobre qué se
cuantifica. De menos a más expresiva:
\termino{proposicional} \(\subset\) \termino{de predicados (primer orden)}
\(\subset\) \termino{de orden superior}.

\smallskip
Cada lógica concreta es una respuesta a las dos. El cálculo del Capítulo 11 es
\emph{clásico y proposicional}; Lean es \emph{clásico (por un axioma) y de
orden superior}.
\end{clave}

\subsection{Las dos preguntas en una tabla}

\begin{figura}{Dónde cae cada sistema de la tesis}
\begin{tikzpicture}[
  celda/.style={caja, text width=3.45cm, minimum height=1.55cm, font=\footnotesize},
  vacia/.style={celda, fill=gray!6, draw=gray!50, text=gray!70!black}]
\node[font=\small\bfseries] at (4.4,1.05) {Mínima};
\node[font=\small\bfseries] at (8.3,1.05) {Intuicionista};
\node[font=\small\bfseries] at (12.2,1.05) {Clásica};
\draw[flecha] (2.6,1.6) -- node[et, above]{más reglas \(\Rightarrow\) más teoremas} (14.1,1.6);
\node[font=\small\bfseries, align=center] at (1.1,0) {Proposicional};
\node[font=\small\bfseries, align=center] at (1.1,-1.8) {De predicados\\(primer orden)};
\node[font=\small\bfseries, align=center] at (1.1,-3.6) {De orden\\superior};
\draw[flecha] (-0.45,0.7) -- node[et, rotate=90, above]{más expresiva} (-0.45,-4.3);
%
\node[vacia] at (4.4,0) {Solo se menciona para aclarar que \(\{\neg,\to\}\) \emph{no} es lógica mínima (Cap.\ 9)};
\node[verde, celda] at (8.3,0) {NJ, BHK, Curry--Howard\\\(\lambda_\to\) (Caps.\ 7 y 9)};
\node[amarillo, celda] at (12.2,0) {\textbf{\(\Provnd\) del Cap.\ 11}\\tablas de verdad\\\(\star\) demostrada en Lean};
\node[vacia] at (4.4,-1.8) {---};
\node[verde, celda] at (8.3,-1.8) {\(\lambda P\): predicados\\como tipos (Cap.\ 8)};
\node[naranja, celda] at (12.2,-1.8) {Completitud de Gödel (1929)\\\emph{fuera de la tesis}};
\node[vacia] at (4.4,-3.6) {---};
\node[verde, celda] at (8.3,-3.6) {CoC y CIC sin axiomas\\(Caps.\ 8 y 10)};
\node[morado, celda] at (12.2,-3.6) {\textbf{Lógica de Lean}\\CIC + \texttt{Classical.choice}\\(nivel meta)};
\end{tikzpicture}

\smallskip
{\footnotesize
\tikz\node[amarillo, minimum width=4mm, minimum height=3mm, inner sep=0pt]{}; demostrado (nivel objeto)\qquad
\tikz\node[morado, minimum width=4mm, minimum height=3mm, inner sep=0pt]{}; lógica en la que se demuestra (nivel meta)\\
\tikz\node[verde, minimum width=4mm, minimum height=3mm, inner sep=0pt]{}; presentado\qquad
\tikz\node[naranja, minimum width=4mm, minimum height=3mm, inner sep=0pt]{}; no demostrado}
\end{figura}

\subsection{Pregunta 1: mínima, intuicionista y clásica}

\begin{figura}{Cada lógica añade una regla a la anterior}
\begin{tikzpicture}
\draw[thick, rounded corners, draw=c3, fill=c3!8] (0,0) rectangle (10.2,5);
\node[anchor=north west, font=\small\bfseries, text=c3] at (0.15,4.95) {Clásica (NK) = intuicionista + RAA};
\draw[thick, rounded corners, draw=c2, fill=c2!12] (0.35,0.35) rectangle (7.1,4.25);
\node[anchor=north west, font=\small\bfseries, text=c2!70!black] at (0.5,4.2) {Intuicionista (NJ) = mínima + \(\bot\)E};
\draw[thick, rounded corners, draw=c1, fill=c1!10] (0.7,0.7) rectangle (3.9,3.45);
\node[align=center, font=\small] at (2.3,2.05) {\textbf{Mínima}\\\(\to\), \(\wedge\), \(\vee\)\\\(\neg\)I, \(\neg\)E\\\footnotesize sin ex falso};
\node[align=center, font=\footnotesize] at (5.5,2.05) {\(\bot\to p\)\\[2pt]\((\neg p\vee q)\to(p\to q)\)\\[2pt]\emph{desde aquí}};
\node[align=center, font=\footnotesize] at (8.65,2.05) {\(p\vee\neg p\)\\[2pt]\(\neg\neg p\to p\)\\[2pt]\(((p\to q)\to p)\to p\)\\[2pt]\emph{solo aquí}};
\begin{scope}[xshift=12.75cm, yshift=2.6cm]
\node[amarillo] (raa) at (90:1.3) {RAA};
\node[amarillo] (lem) at (210:1.3) {LEM};
\node[amarillo] (dne) at (330:1.3) {\(\neg\neg\)E};
\draw[dflecha] (raa) -- (lem); \draw[dflecha] (lem) -- (dne); \draw[dflecha] (dne) -- (raa);
\node[font=\scriptsize, align=center] at (0,-2.0) {cualquiera de las tres\\añadida a NJ da\\la lógica clásica};
\end{scope}
\end{tikzpicture}
\end{figura}

\paragraph{Lógica clásica.} Cada fórmula es verdadera o falsa, sin tercera
opción (\termino{bivalencia}); su semántica son las tablas de verdad. La
caracteriza cualquiera de tres principios equivalentes: el tercero excluido
(LEM, \(\varphi\vee\neg\varphi\)), la eliminación de la doble negación
(\(\neg\neg\varphi\Rightarrow\varphi\)) y la reducción al absurdo clásica (RAA).
El cálculo \(\Provnd\) del Capítulo 11 es clásico; su única regla clásica es RAA,
y dentro del cálculo solo se usa en el Lema 11.11.

\paragraph{Lógica intuicionista} (también \termino{lógica de Heyting} o
constructiva). Ser verdadero es \emph{tener una construcción}: no basta mostrar
que algo no puede ser falso. La interpretación BHK (Cap.\ 9) dice qué es una
prueba de cada conectivo; Curry--Howard la convierte en programas.

\begin{center}
\renewcommand{\arraystretch}{1.25}
{\small
\begin{tabular}{@{}l >{\raggedright\arraybackslash}p{6.3cm} >{\raggedright\arraybackslash}p{4.6cm}@{}}
\toprule
\textbf{Fórmula} & \textbf{Una prueba es\dots{} (BHK)} & \textbf{En Curry--Howard / Lean}\\
\midrule
\(\varphi\wedge\psi\) & un par: una prueba de \(\varphi\) y una de \(\psi\) & par \(\langle a,b\rangle\)\\
\(\varphi\vee\psi\) & una prueba de \emph{uno} de los dos, \emph{diciendo cuál} & \texttt{Or.inl a} o \texttt{Or.inr b}\\
\(\varphi\to\psi\) & un método que convierte pruebas de \(\varphi\) en pruebas de \(\psi\) & función \(\lambda x.\,M\)\\
\(\bot\) & no hay ninguna & tipo vacío \texttt{False}\\
\(\neg\varphi\) (\(=\varphi\to\bot\)) & un método que convierte pruebas de \(\varphi\) en absurdo & función hacia \texttt{False}\\
\(\forall x.\,P(x)\)\(^{*}\) & un método que a cada \(x\) le asigna una prueba de \(P(x)\) & función dependiente (\(\Pi\))\\
\(\exists x.\,P(x)\)\(^{*}\) & un testigo \(a\) junto con una prueba de \(P(a)\) & par dependiente\\
\bottomrule
\end{tabular}}

{\footnotesize \(^{*}\)Extensión estándar; la definición de la tesis cubre \(\wedge\), \(\vee\), \(\to\) y \(\bot\).}
\end{center}

\textbf{Por qué falla \(p\vee\neg p\).} BHK exige decir \emph{cuál} de los dos
lados vale, y en general no se sabe. Ejemplo estándar (fuera de la tesis):
existen irracionales \(a,b\) con \(a^b\) racional. Si \(\sqrt2^{\sqrt2}\) es
racional, sirven \(a=b=\sqrt2\); si no, sirven \(a=\sqrt2^{\sqrt2}\) y
\(b=\sqrt2\), porque \(a^b=\sqrt2^{\,2}=2\). La prueba es clásicamente
correcta, pero no dice cuál de los dos pares funciona: no es constructiva.

\textbf{Lo que se conserva.} El ex falso (\(\bot\)E) y la introducción de la
negación: para probar \(\neg A\) se supone \(A\) y se llega a \(\bot\). Eso
\emph{no} es RAA; es \(\to\)I, porque \(\neg A\) abrevia \(A\to\bot\).

\begin{center}
\begin{tabular}{c@{\qquad}c}
\AxiomC{\(\Gamma,A\vdash\bot\)}\RightLabel{\scriptsize\(\neg\)I (intuicionista)}\UnaryInfC{\(\Gamma\vdash\neg A\)}\DisplayProof
&
\AxiomC{\(\Gamma,\neg A\vdash\bot\)}\RightLabel{\scriptsize RAA (clásica)}\UnaryInfC{\(\Gamma\vdash A\)}\DisplayProof
\end{tabular}
\end{center}

\paragraph{Lógica mínima.} Es la intuicionista \emph{sin} ex falso: \(\bot\)
es un átomo más, sin poderes especiales. La tesis la menciona una sola vez
(Cap.\ 9), para aclarar que usar la base \(\{\neg,\to\}\) \emph{no} es lógica
mínima: es lógica clásica escrita con menos conectivos primitivos.

\newcommand{\lsi}{\textcolor{c2!80!black}{\textbf{sí}}}
\newcommand{\lno}{\textcolor{c3}{no}}

\begin{center}
\renewcommand{\arraystretch}{1.2}
{\small
\begin{tabular}{@{}l c c c >{\raggedright\arraybackslash}p{5.2cm}@{}}
\toprule
\textbf{Fórmula} & \textbf{Mín.} & \textbf{Intuic.} & \textbf{Clás.} & \textbf{Comentario}\\
\midrule
\(p\to\neg\neg p\) & \lsi & \lsi & \lsi & la doble negación se \emph{introduce} siempre\\
\(\neg\neg\neg p\to\neg p\) & \lsi & \lsi & \lsi & tres negaciones se reducen a una\\
\(\neg\neg(p\vee\neg p)\) & \lsi & \lsi & \lsi & LEM doblemente negado\\
\((\neg p\vee\neg q)\to\neg(p\wedge q)\) & \lsi & \lsi & \lsi & De Morgan, dirección fácil\\
\midrule
\(\bot\to p\) & \lno & \lsi & \lsi & ex falso (\(\bot\)E)\\
\((\neg p\vee q)\to(p\to q)\) & \lno & \lsi & \lsi & el caso \(\neg p\) necesita ex falso\\
\midrule
\(p\vee\neg p\) & \lno & \lno & \lsi & tercero excluido (LEM)\\
\(\neg\neg p\to p\) & \lno & \lno & \lsi & eliminación de la doble negación\\
\(((p\to q)\to p)\to p\) & \lno & \lno & \lsi & ley de Peirce\\
\(\neg(p\wedge q)\to(\neg p\vee\neg q)\) & \lno & \lno & \lsi & De Morgan, dirección difícil\\
\((p\to q)\to(\neg p\vee q)\) & \lno & \lno & \lsi & \(\to\) como disyunción\\
\bottomrule
\end{tabular}}
\end{center}

\begin{nota}[Resultado estándar, fuera de la tesis: teorema de Glivenko]
Una fórmula proposicional \(\varphi\) es tautología clásica si y solo si
\(\neg\neg\varphi\) es teorema intuicionista. La lógica intuicionista no
pierde verdades clásicas: las conserva «bajo doble negación».
\end{nota}

\subsection{La semántica de la lógica intuicionista: modelos de Kripke}

Las tablas de verdad no sirven para la lógica intuicionista: NJ no prueba
\(p\vee\neg p\), que es tautología. Su semántica son los \termino{modelos de
Kripke}: \emph{mundos} (estados de conocimiento) ordenados por «saber más».
Lo que se sabe en un mundo se sigue sabiendo en los posteriores. La tesis los
cita (Cap.\ 9, Sørensen--Urzyczyn) pero no los desarrolla; el anexo de Kripke
está excluido de la versión compilada.

\begin{figura}{El contramodelo más pequeño de \(p\vee\neg p\)}
\begin{tikzpicture}
\node[circle, draw=c1, fill=c1!10, minimum size=1.7cm, align=center] (w0) at (0,0) {\(w_0\)\\\footnotesize \(p\) aún\\\footnotesize no se sabe};
\node[circle, draw=c2, fill=c2!14, minimum size=1.7cm, align=center] (w1) at (4.2,0) {\(w_1\)\\\footnotesize \(p\) verdadera};
\draw[flecha] (w0) -- node[et, above]{se aprende más} (w1);
\node[align=left, font=\small, anchor=west] at (6.2,0) {%
\(w_0\nVdash p\)\\
\(w_0\nVdash\neg p\), porque en \(w_1\) llega a valer \(p\)\\
luego \(w_0\nVdash p\vee\neg p\)\\[3pt]
además \(w_0\Vdash\neg\neg p\) pero \(w_0\nVdash p\):\\
falla \(\neg\neg p\to p\)};
\end{tikzpicture}
\end{figura}

\begin{center}
{\small
\renewcommand{\arraystretch}{1.2}
\begin{tabular}{@{}l@{\qquad}l@{}}
\toprule
\textbf{Cuándo} & \textbf{significa}\\
\midrule
\(w\Vdash\varphi\wedge\psi\) & \(w\Vdash\varphi\) y \(w\Vdash\psi\)\\
\(w\Vdash\varphi\vee\psi\) & \(w\Vdash\varphi\) o \(w\Vdash\psi\), \emph{ya} en \(w\)\\
\(w\Vdash\varphi\to\psi\) & en todo \(w'\geq w\): si \(w'\Vdash\varphi\), entonces \(w'\Vdash\psi\)\\
\(w\Vdash\neg\varphi\) & ningún \(w'\geq w\) cumple \(w'\Vdash\varphi\)\\
\bottomrule
\end{tabular}}
\end{center}

Se lee «\(w\) \termino{fuerza} \(\varphi\)». Con un solo mundo, las cláusulas
coinciden con las tablas de verdad: la lógica clásica es el caso en que ya no
se aprende nada nuevo.

\subsection{Pregunta 2: proposicional, de predicados y de orden superior}

\begin{center}
\renewcommand{\arraystretch}{1.3}
{\small
\begin{tabular}{@{}>{\raggedright\arraybackslash}p{3.0cm} >{\raggedright\arraybackslash}p{4.4cm} >{\raggedright\arraybackslash}p{3.9cm} >{\raggedright\arraybackslash}p{3.5cm}@{}}
\toprule
\textbf{Lógica} & \textbf{Qué añade} & \textbf{Ejemplo} & \textbf{En la tesis}\\
\midrule
proposicional & átomos \(p,q,\dots\) y conectivos; sin cuantificadores & \(p\to(q\to p)\) & Caps.\ 9 y 11; \(\lambda_\to\)\\
de predicados (primer orden) & predicados sobre individuos y \(\forall x\), \(\exists x\) sobre individuos & \(\forall x\,(P(x)\to\exists y\,R(x,y))\) & \(\lambda P\) (Cap.\ 8); Frege; su completitud no se demuestra\\
de orden superior & cuantificar sobre predicados y proposiciones & \(\forall P.\,(P\to P)\), con \(P\) una proposición & CoC y CIC (Caps.\ 8 y 10); Lean\\
\bottomrule
\end{tabular}}
\end{center}

El ejemplo más importante de orden superior es el \emph{principio de
inducción}, que cuantifica sobre todos los predicados \(P\):
\[
\forall P.\;\bigl(P(0)\wedge\forall n\,(P(n)\to P(n+1))\bigr)\to\forall n\,P(n).
\] Por eso en CIC la inducción es un único
teorema (el eliminador) y no un esquema con una instancia por fórmula.

\subsection{Las dos lógicas de la demostración del Capítulo 11}

\begin{figura}{El cálculo estudiado vive dentro de la lógica que lo estudia}
\begin{tikzpicture}
\draw[thick, rounded corners, draw=c4, fill=c4!8] (0,0) rectangle (15,3.9);
\node[anchor=north west, font=\small\bfseries, text=c4!80!black, align=left] at (0.15,3.85) {Nivel meta: la lógica de Lean = CIC + \texttt{propext} + \texttt{Classical.choice} + \texttt{Quot.sound}\\\footnotesize\mdseries clásica (por \texttt{Classical.choice}) y de orden superior};
\draw[thick, rounded corners, draw=c5, fill=c5!15] (0.4,0.35) rectangle (7.4,2.55);
\node[align=center, font=\small] at (3.9,1.45) {\textbf{Nivel objeto: el cálculo \(\Provnd\)}\\clásico, proposicional, base \(\{\neg,\to\}\)\\\footnotesize RAA es un \emph{constructor} del predicado inductivo};
\node[align=left, font=\small, anchor=west] at (7.8,1.45) {Lean demuestra teoremas \emph{sobre} \(\Provnd\):\\\(\mathrm{Taut}(\varphi)\Rightarrow\varnothing\Provnd\varphi\).\\[2pt]\footnotesize El análisis por casos sobre si \(v(p)\) es verdadera\\\footnotesize usa \texttt{Classical.em}, derivado de \texttt{Classical.choice}.};
\end{tikzpicture}
\end{figura}

\begin{nota}[Pregunta probable: «Si Curry--Howard es intuicionista, ¿cómo vive en Lean una prueba clásica?»]
Hay que separar los dos niveles. \emph{En el objeto}, RAA es solo un
constructor del predicado inductivo \(\Provnd\): construir una derivación que lo
use no requiere nada clásico de Lean. \emph{En el meta}, la demostración
razona por casos sobre el valor de verdad de los átomos; como las valuaciones
viven en \texttt{Prop}, ese paso usa \texttt{Classical.choice}, y la auditoría
de axiomas lo confirma. El núcleo de Lean (CIC) es intuicionista; el axioma lo
vuelve clásico. Más detalle en la sección «Nivel objeto y nivel meta».
\end{nota}

\subsection{Dónde aparece cada lógica}

\begin{center}
\renewcommand{\arraystretch}{1.25}
{\small
\begin{tabular}{@{}>{\raggedright\arraybackslash}p{3.6cm} >{\raggedright\arraybackslash}p{5.6cm} >{\raggedright\arraybackslash}p{5.6cm}@{}}
\toprule
\textbf{Lógica} & \textbf{Dónde aparece} & \textbf{Qué se hace con ella}\\
\midrule
clásica proposicional & Cap.\ 9 (NK, tablas), Cap.\ 11, Ap.\ A & \textbf{demostrada en Lean}: completitud débil, corrección, adecuación\\
intuicionista proposicional (NJ) & Cap.\ 9: reglas, BHK, Curry--Howard & presentada; no formalizada\\
semántica de Kripke & Cap.\ 9 (nota); anexo excluido & citada (Sørensen--Urzyczyn)\\
mínima & Cap.\ 9, una mención & solo para aclarar que \(\{\neg,\to\}\) no es lógica mínima\\
de predicados / primer orden & contexto histórico (Frege), Cap.\ 8 (\(\lambda P\)), conclusiones & presentada como sistema de tipos; su completitud no se demuestra\\
intuicionista de orden superior & Cap.\ 10 (Coquand--Huet) & la lógica para la que se diseñó el CoC\\
lógica de Lean & glosario de la tesis, auditorías de axiomas & la lógica en la que se demuestra todo\\
de Heyting & Cap.\ 10, diagrama & otro nombre de la intuicionista\\
\bottomrule
\end{tabular}}
\end{center}

\begin{nota}[No confundir]
\begin{itemize}
\item \emph{Lógica mínima} (sin ex falso) \(\neq\) \emph{base mínima de
  conectivos} (\(\{\neg,\to\}\), que es lógica clásica completa).
\item \(\neg\)I (suponer \(A\), llegar a \(\bot\), concluir \(\neg A\)) es
  intuicionista; RAA (suponer \(\neg A\), llegar a \(\bot\), concluir \(A\)) es
  clásica.
\item «Intuicionista» no significa «intuitiva»: viene del intuicionismo de
  Brouwer, que exige construcciones.
\item «Lógica inclusiva» no es una lógica: la tesis dice «disyunción lógica
  inclusiva» para el \(\vee\) que vale si vale uno de los lados o ambos.
\end{itemize}
\end{nota}

%=====================================================================
\section{Bases de conectivos y completitud funcional}
%=====================================================================

\begin{figura}{Una fórmula define una función de verdad}
\begin{tikzpicture}
\node[azul, text width=34mm] (in) at (0,0) {\(\{\mathsf V,\mathsf F\}^n\)\\filas de la tabla\\(\(2^n\) filas)};
\node[verde, text width=30mm] (out) at (7,0) {\(\{\mathsf V,\mathsf F\}\)\\valor de la fórmula};
\draw[flecha] (in) -- node[et,above]{función de verdad \(f\)} (out);
\node[font=\small, align=center] at (3.5,-1.4) {hay \(2^{2^n}\) funciones de verdad de \(n\) argumentos.\\Una \textbf{base} es \textbf{funcionalmente completa} si con ella se escribe \emph{cualquiera} de ellas.};
\end{tikzpicture}
\end{figura}

\begin{figura}{Cómo se ve que \(\{\neg,\to\}\) basta}
\begin{tikzpicture}
\node[azul, text width=40mm] (t) at (0,0) {\textbf{1.} Tabla cualquiera\\\(\Rightarrow\) forma normal disyuntiva:\\una conjunción de literales\\por cada fila verdadera};
\node[verde, text width=40mm] (n) at (5.4,0) {\textbf{2.} \(\wedge,\vee\) se escriben con \(\neg,\to\):\\\(A\vee B\equiv\neg A\to B\)\\\(A\wedge B\equiv\neg(A\to\neg B)\)};
\node[amarillo, text width=36mm] (r) at (10.4,0) {\textbf{3.} Toda función de verdad se escribe con \(\neg\) y \(\to\)};
\draw[flecha] (t) -- (n); \draw[flecha] (n) -- (r);
\end{tikzpicture}
\end{figura}

\begin{center}
{\small
\begin{tabular}{@{}lll@{}}
\toprule
\textbf{Base} & \textbf{¿Completa?} & \textbf{Por qué}\\
\midrule
\(\{\neg,\to\}\), \(\{\neg,\wedge\}\), \(\{\neg,\vee\}\) & sí & expresan \(\wedge,\vee,\neg\)\\
\(\{\mathsf{NAND}\}\), \(\{\mathsf{NOR}\}\) & sí & un solo conectivo basta\\
\(\{\wedge,\vee\}\), \(\{\to\}\), \(\{\wedge,\vee,\to\}\) & no & con todo \(\mathsf V\) dan \(\mathsf V\): no expresan \(\neg p\)\\
\bottomrule
\end{tabular}}
\end{center}

\begin{nota}[No confundir dos «completitudes»]
\textbf{Completitud funcional} habla de lo que se puede \emph{expresar} (una
propiedad del lenguaje). \textbf{Completitud del cálculo} habla de lo que se
puede \emph{demostrar} (una propiedad de las reglas). La tesis usa la primera
para justificar la base \(\{\neg,\to\}\) y demuestra la segunda.
\end{nota}

%=====================================================================
\section{El Cálculo de Construcciones Inductivas}
%=====================================================================

\begin{figura}{Anatomía de un tipo inductivo: cuatro piezas de una sola declaración}
\begin{tikzpicture}
\node[amarillo, text width=30mm] (f) at (0,0) {\textbf{Formación}\\\(\mathbb N\) es un tipo};
\node[azul, text width=30mm] (i) at (3.9,0) {\textbf{Introducción}\\constructores \(0\), \(S\)\\(únicas formas de construir)};
\node[verde, text width=30mm] (e) at (7.8,0) {\textbf{Eliminación}\\recursor / inducción\\(única forma de usar)};
\node[morado, text width=30mm] (c) at (11.7,0) {\textbf{Cómputo}\\reducción \(\iota\):\\\(R(0)\to x\), \(R(Sn)\to f\,n\,R(n)\)};
\draw[flecha] (f) -- (i); \draw[flecha] (i) -- (e); \draw[flecha] (e) -- (c);
\end{tikzpicture}
\end{figura}

\begin{figura}{Dos maneras de tener números naturales}
\begin{tikzpicture}
\node[naranja, text width=62mm] (coc) at (0,0) {\textbf{CoC: codificación de Church}\\\(n:=\lambda A\,x\,f.\,f^{\,n}x\)\\«\(n\) es aplicar \(f\) \(n\) veces»\\[2pt]\checkmark\ iteración\quad \(\times\) inducción derivable\\\(\times\) predecesor caro};
\node[verde, text width=62mm] (cic) at (8,0) {\textbf{CIC: tipo inductivo}\\\(0:\mathbb N,\quad S:\mathbb N\to\mathbb N\)\\el eliminador \emph{es} la inducción\\[2pt]\checkmark\ recursión\quad \checkmark\ inducción\\\checkmark\ cómputo por \(\iota\)};
\draw[flecha] (coc) -- node[et,above]{CIC} (cic);
\end{tikzpicture}
\end{figura}

%=====================================================================
\section{Capítulo 11 en un diagrama}
%=====================================================================

\begin{figura}{De la tabla de verdad a una derivación sin hipótesis}
\begin{tikzpicture}[font=\small]
\node[azul, text width=34mm] (tab) at (0,0) {\textbf{Tabla de verdad}\\\(\varphi\) verdadera en\\las \(2^k\) filas};
\node[azul, text width=38mm] (lit) at (4.8,0) {\textbf{Contexto de literales}\\fila \(v\): \(\Gamma_v=\{p\text{ ó }\neg p\}\)\\ej.\ \(\{\neg q,\,p\}\)};
\node[verde, text width=36mm] (kal) at (9.6,0) {\textbf{Kalmár (11.15)}\\\(\Gamma_v\vdash\varphi\)\\en cada fila};
\node[naranja, text width=40mm] (des) at (9.6,-2.6) {\textbf{Descarga (11.17)}\\juntar filas que difieren\\en un átomo (11.11)};
\node[amarillo, text width=34mm] (fin) at (4.8,-2.6) {\(\varnothing\vdash\varphi\)};
\draw[flecha] (tab) -- (lit); \draw[flecha] (lit) -- (kal); \draw[flecha] (kal) -- (des); \draw[flecha] (des) -- (fin);
\end{tikzpicture}
\end{figura}

\begin{figura}{El árbol de descarga con dos átomos}
\begin{tikzpicture}[level distance=13mm, level 1/.style={sibling distance=60mm}, level 2/.style={sibling distance=30mm},
  every node/.style={caja, fill=c1!8, draw=c1}, lab/.style={draw=none, fill=none, font=\scriptsize}]
\node[fill=c5!20] {\(\varnothing\vdash\varphi\)}
  child { node {\(\{p\}\vdash\varphi\)}
    child { node[fill=c2!15] {\(\{q,p\}\vdash\varphi\)} edge from parent node[lab,left]{\(q\)} }
    child { node[fill=c2!15] {\(\{\neg q,p\}\vdash\varphi\)} edge from parent node[lab,right]{\(\neg q\)} }
    edge from parent node[lab,left]{\(p\)} }
  child { node {\(\{\neg p\}\vdash\varphi\)}
    child { node[fill=c2!15] {\(\{q,\neg p\}\vdash\varphi\)} edge from parent node[lab,left]{\(q\)} }
    child { node[fill=c2!15] {\(\{\neg q,\neg p\}\vdash\varphi\)} edge from parent node[lab,right]{\(\neg q\)} }
    edge from parent node[lab,right]{\(\neg p\)} };
\end{tikzpicture}
{\small hojas: Kalmár; cada nodo interno: análisis por casos (usa RAA).}
\end{figura}

%=====================================================================
\section{Capítulo 12 en un diagrama}
%=====================================================================

\begin{figura}{Partición, recursión y ensamblaje}
\begin{tikzpicture}[font=\small]
\node[caja, fill=c1!10] (p) at (0,0) {\(p\)};
\node[caja, fill=gray!10, right=1mm of p, minimum width=60mm] (r) {cola \(r\) sin ordenar};
\node[verde] (L) at (1.5,-1.6) {\(L\): los \(\leq p\)};
\node[naranja] (R) at (6.5,-1.6) {\(R\): los \(>p\)};
\node[verde] (A) at (1.5,-3.2) {\(\qs(L)\) ordenada};
\node[caja, fill=c1!10] (pp) at (4,-3.2) {\(p\)};
\node[naranja] (B) at (6.5,-3.2) {\(\qs(R)\) ordenada};
\draw[flecha] (r.south) -- ++(0,-0.3) -| (L.north);
\draw[flecha] (r.south) -- ++(0,-0.3) -| (R.north);
\draw[flecha] (L) -- node[et,left]{recursión} (A);
\draw[flecha] (R) -- node[et,right]{recursión} (B);
\node[draw, dashed, rounded corners, fit=(A)(pp)(B), inner sep=4pt, label=below:{\(\qs(p::r)=\qs(L)\concat(p::\qs(R))\)}] {};
\node[font=\scriptsize, align=left, anchor=west] at (9,-1) {\textbf{terminación}: \(|L|,|R|<|p::r|\)\\\textbf{permutación}: \(L\concat R\sim r\)\\\textbf{orden}: \(\qs(L)\leq p\leq\qs(R)\)};
\end{tikzpicture}
\end{figura}

\begin{figura}{La especificación: dos condiciones, ninguna basta sola}
\begin{tikzpicture}
\node[verde, text width=44mm] (o) at (0,0) {\textbf{Ordenada}\\\(f(\ell)=[\,]\) la cumple\\pero pierde elementos};
\node[azul, text width=44mm] (p) at (6,0) {\textbf{Permutación}\\\(f(\ell)=\ell\) la cumple\\pero no ordena};
\node[amarillo, text width=44mm] (c) at (3,-2.2) {\textbf{Correcta} = las dos a la vez};
\draw[flecha] (o) -- (c); \draw[flecha] (p) -- (c);
\end{tikzpicture}
\end{figura}

%=====================================================================
\section{¿Qué significa «verificado»?}
%=====================================================================

\begin{figura}{La cadena de confianza}
\begin{tikzpicture}[font=\small]
\node[azul, text width=28mm] (a) at (0,0) {enunciado\\matemático};
\node[azul, text width=28mm] (b) at (3.7,0) {definición\\formal};
\node[verde, text width=28mm] (c) at (7.4,0) {prueba\\(término)};
\node[morado, text width=28mm] (d) at (11.1,0) {núcleo\\comprueba el tipo};
\node[amarillo, text width=28mm] (e) at (11.1,-2.3) {auditoría:\\¿qué axiomas?};
\draw[flecha] (a) -- node[et,above]{\(\star\)} (b); \draw[flecha] (b) -- (c); \draw[flecha] (c) -- (d); \draw[flecha] (d) -- (e);
\node[naranja, text width=80mm, font=\small] at (3.7,-2.3) {\(\star\) El único paso que la máquina \emph{no} comprueba: que la definición formal dice lo que se quería decir. Por eso hay que defender las especificaciones.};
\end{tikzpicture}
\end{figura}

%=====================================================================
\section{Glosario}
%=====================================================================

Los términos van en orden alfabético. Entre corchetes, dónde aparecen.

{\small
\begin{longtable}{@{}p{0.27\textwidth}p{0.69\textwidth}@{}}
\toprule
\textbf{Término} & \textbf{Significado}\\
\midrule
\endhead
abstracción & Construir una función: \(\lambda x.M\) es «la función que a \(x\) le asigna \(M\)». [Cap.\ 7]\\
adecuación & Corrección + completitud: \(\mathrm{Taut}(\varphi)\iff\varnothing\vdash\varphi\). [Ap.\ A]\\
\(\alpha\)-equivalencia & Igualdad salvo nombres de variables ligadas: \(\lambda x.x=_\alpha\lambda y.y\). [7.2.4]\\
alcance & Parte del término donde una variable ligada está disponible; en las pruebas, la «caja» de una suposición.\\
antecedente / consecuente & En \(\varphi\to\psi\), \(\varphi\) es el antecedente y \(\psi\) el consecuente.\\
antisimétrica & \(a\leq b\) y \(b\leq a\) implican \(a=b\).\\
aplicación & Usar una función: \(M\,N\) es «\(M\) aplicada a \(N\)». [Cap.\ 7]\\
aridad & Número de argumentos de una función o conectivo; en el CIC, la forma de un tipo inductivo. [10.3.2]\\
asistente de pruebas & Programa que comprueba demostraciones formales con un núcleo pequeño.\\
átomo & Variable proposicional: \(p,q,r\). Fórmula sin conectivos.\\
auditoría de axiomas & Lista de los axiomas de los que depende un teorema; detecta huecos (\texttt{sorry}). [11.6, 12.10]\\
axioma & Enunciado que se acepta sin demostración. En el asistente: extensionalidad proposicional, elección y cocientes.\\
\(\beta\)-reducción & Paso de cómputo: \((\lambda x.M)N\bred M[x:=N]\). [7.2.7]\\
base (de conectivos) & Conjunto de conectivos primitivos del lenguaje. La tesis usa \(\{\neg,\to\}\). [11.2.2]\\
BHK & Interpretación de Brouwer--Heyting--Kolmogorov: una prueba es una construcción. [9.4.1]\\
bien fundada (relación) & Sin cadenas descendentes infinitas; permite la inducción bien fundada. [12.4]\\
bivalencia & Toda proposición es verdadera o falsa. Hace válida la RAA en tablas de verdad.\\
caja (de suposición) & Zona de una prueba donde vale una suposición; se cierra al descargarla. [Nota 9.5]\\
captura & Error en el que una variable libre queda ligada al sustituir. [Def.\ 7.17]\\
CIC & Cálculo de Construcciones Inductivas: CoC + tipos inductivos. Base del asistente. [Cap.\ 10]\\
clausura reflexiva-transitiva & Menor relación reflexiva y transitiva que contiene a otra; \(\bmred\) es la de \(\bred\).\\
CoC & Cálculo de Construcciones, el vértice \(\lambda C\) del cubo. [8.4]\\
codificación de Church & Representar datos como funciones: \(n=\lambda A\,x\,f.\,f^nx\). [10.2]\\
cobertura & Lista \(A\) con \(\mathrm{atomos}(\varphi)\subseteq A\). [11.5.2]\\
completitud (del cálculo) & Lo verdadero es demostrable: \(\models\varphi\Rightarrow\vdash\varphi\). [Cap.\ 11]\\
completitud funcional & Una base puede expresar toda función de verdad. [11.2.2]\\
concatenación & Pegar dos listas: \([1,2]\concat[3]=[1,2,3]\).\\
confluencia & Si un término se reduce a dos, ambos se reducen a uno común. [Teo.\ 7.37]\\
congruencia & Relación compatible con los constructores: si \(M\,R\,N\), entonces \(\lambda x.M\,R\,\lambda x.N\), etc.\\
conectivo & Símbolo que combina fórmulas: \(\neg,\wedge,\vee,\to\).\\
consecuencia semántica & \(\Gamma\models\varphi\): toda valuación que satisface \(\Gamma\) hace verdadera \(\varphi\). [9.2]\\
consistencia & No se deriva \(\bot\). [Cor.\ A.5]\\
constructiva (prueba) & Prueba que exhibe lo que afirma: un testigo para \(\exists\), el lado concreto para \(\vee\). Es el criterio de la lógica intuicionista. [9.4]\\
constructor & Regla que genera elementos de un tipo inductivo: \(0\), \(S\), \([\,]\), \(::\).\\
contexto & Conjunto (o lista) de hipótesis o declaraciones: \(\Gamma\).\\
contractum & Resultado de contraer un redex: \(M[x:=N]\).\\
contradicción & Fórmula falsa en todas las filas; o, en una prueba, obtener \(\bot\).\\
corrección (de un cálculo) & Lo demostrable es verdadero (solidez). [Teo.\ A.3]\\
corrección (de un algoritmo) & El algoritmo cumple su especificación. [Cap.\ 12]\\
Curry--Howard & Proposiciones = tipos; pruebas = programas. [9.5]\\
debilitamiento & Añadir hipótesis no invalida una derivación. [Lema 11.9]\\
decidible & Existe un algoritmo que siempre termina y responde sí o no.\\
deducción natural & Sistema de pruebas con reglas de introducción y eliminación por conectivo. [Cap.\ 9]\\
derivación & Árbol de reglas que termina en un juicio; una prueba formal.\\
descarga & Cerrar una suposición al aplicar una regla (\(\to\)I, \(\neg\)I, RAA, \(\vee\)E).\\
diamante (propiedad) & Si \(M\) va a \(N_1\) y a \(N_2\), existe \(L\) al que van ambos. [7.2.8]\\
divide y vencerás & Partir el problema, resolver las partes recursivamente y combinar. [12.3]\\
eliminador / recursor & Única forma de usar un tipo inductivo; da recursión e inducción. [10.3.4]\\
especificación & Enunciado formal de lo que debe cumplir un programa.\\
explosión (ex falso) & De \(\bot\) se deriva cualquier fórmula (\(\bot\)E).\\
fórmula & Expresión bien formada del lenguaje lógico.\\
forma normal (\(\beta\)) & Término sin redexes. [Def.\ 7.33]\\
forma normal disyuntiva & Disyunción de conjunciones de literales; existe para toda tabla.\\
función total / parcial & Total: definida en todo argumento. Parcial: puede no estarlo (no terminar).\\
generalizar (la hipótesis) & Cuantificar una variable antes de inducir para que la hipótesis inductiva sirva en otros casos.\\
Glivenko (teorema de) & \(\varphi\) es tautología clásica si y solo si \(\neg\neg\varphi\) es teorema intuicionista. Resultado estándar, fuera de la tesis.\\
Heyting (lógica de) & Otro nombre de la lógica intuicionista. [Cap.\ 10]\\
hipótesis inductiva & Lo que se supone para los casos menores en una inducción.\\
impredicativo & Definición que cuantifica sobre una colección que incluye lo que se define. [10.2]\\
inducción bien fundada & Suponer la propiedad para todo lo «más pequeño». [12.4]\\
inducción estructural & Un caso por constructor. [Nota 7.5]\\
inducción sobre derivaciones & Un caso por regla. [Lemas 7.57, 11.9]\\
juicio & Afirmación del sistema formal: \(\Gamma\vdash M:\sigma\) o \(\Gamma\vdash\varphi\).\\
Kalmár (lema de) & En cada fila, los literales derivan \(\varphi\) o \(\neg\varphi\). [Lema 11.15]\\
kind & Tipo de los tipos: \(*\), \(*\to*\). [8.2]\\
lenguaje objeto / metalenguaje & El sistema estudiado frente al lenguaje en que se razona sobre él.\\
LEM & Tercero excluido: \(\varphi\vee\neg\varphi\).\\
literal & Un átomo o su negación: \(p\) o \(\neg p\). [Def.\ 11.13]\\
lógica clásica & Bivalente: incluye LEM, \(\neg\neg\)E o RAA (equivalentes). Su semántica son las tablas de verdad. [Caps.\ 9 y 11]\\
lógica de Lean & CIC con \texttt{propext}, \texttt{Classical.choice} y \texttt{Quot.sound}. Es clásica porque \texttt{Classical.choice} implica LEM.\\
lógica de orden superior & Permite cuantificar sobre predicados y proposiciones: \(\forall P.\,P\to P\). [Cap.\ 10]\\
lógica de predicados (primer orden) & Añade predicados y cuantificadores sobre individuos. Su completitud (Gödel, 1929) no está en la tesis. [Cap.\ 8, \(\lambda P\)]\\
lógica intuicionista & Una prueba es una construcción (BHK). Sin LEM, \(\neg\neg\)E ni RAA; con ex falso. Su cálculo es NJ. [9.4]\\
lógica mínima & La intuicionista sin ex falso (\(\bot\)E). No es la base \(\{\neg,\to\}\). [Cap.\ 9]\\
lógica proposicional & Solo átomos y conectivos, sin cuantificadores. [Caps.\ 9 y 11]\\
medida & Función hacia \(\mathbb N\) que decrece en cada llamada recursiva: \(|\ell|\). [12.4]\\
metateorema & Teorema \emph{sobre} un sistema formal (corrección, completitud).\\
modelo de Kripke & Semántica de la lógica intuicionista: mundos ordenados por «saber más»; lo que vale en un mundo sigue valiendo en los posteriores. [Cap.\ 9, citado]\\
modus ponens & De \(\varphi\) y \(\varphi\to\psi\), concluir \(\psi\) (\(\to\)E).\\
multiplicidad & Número de veces que un elemento aparece en una lista. [12.2]\\
NJ / NK & Deducción natural intuicionista / clásica; NK = NJ + RAA. [Cap.\ 9]\\
normalización débil / fuerte & Alguna / toda ruta de reducción termina. [Def.\ 7.35]\\
núcleo (kernel) & Parte mínima del asistente que comprueba los tipos.\\
orden lineal & Reflexivo, antisimétrico, transitivo y total. [12.2]\\
PAT & «Propositions as types»: la lectura de Curry--Howard. [Def.\ 7.61]\\
Peirce (ley de) & \(((p\to q)\to p)\to p\): tautología clásica que no es teorema intuicionista.\\
permutación & Misma lista salvo el orden (mismas multiplicidades). [Def.\ 12.3]\\
pivote & Elemento que separa la lista en Quicksort.\\
positividad estricta & La variable recursiva no aparece a la izquierda de \(\to\) en un constructor. [10.3.2]\\
preorden & Relación reflexiva y transitiva.\\
\(\Pi\) (producto dependiente) & Tipo de funciones cuyo resultado depende del argumento; codifica \(\forall\). [8.1.1]\\
RAA & Reducción al absurdo clásica: de \(\Gamma,\neg\varphi\vdash\bot\), concluir \(\varphi\).\\
redex & Subtérmino reducible \((\lambda x.M)N\).\\
reducción \(\iota\) & El eliminador aplicado a un constructor se calcula con su caso. [10.3.5]\\
regla derivada & Regla que se demuestra a partir de las primitivas (análisis por casos, Lema 11.11).\\
satisfacible & Verdadera en al menos una fila.\\
secuente & Expresión \(\Gamma\vdash\varphi\). [Def.\ 9.3]\\
sort & \(*\) o \(\square\): las clasificaciones del nivel superior. [Cap.\ 8]\\
subtérmino & Parte de un término: \(\mathrm{Sub}(M)\). [Def.\ 7.3]\\
sustitución & \(M[x:=N]\): reemplazar las \(x\) libres por \(N\), evitando captura. [Def.\ 7.20]\\
tautología & Verdadera en todas las filas. [Def.\ 11.5]\\
teorema (del cálculo) & Fórmula derivable sin hipótesis: \(\vdash\varphi\). [Def.\ 9.4]\\
terminación & Que toda ejecución acabe. [Prop.\ 12.6]\\
tipo & Clasificación de un término; bajo PAT, una proposición.\\
tipo inductivo & Tipo definido por constructores. [Cap.\ 10]\\
universo & Tipo de tipos en una jerarquía (\(\mathrm{Type}_0:\mathrm{Type}_1:\cdots\)).\\
valuación & Asignación de \(\mathsf V\) o \(\mathsf F\) a cada átomo; una fila de la tabla. [Def.\ 11.4]\\
variable libre / ligada & Libre: no está bajo un \(\lambda\) que la enlace. Ligada: sí lo está.\\
\bottomrule
\end{longtable}
}

%=====================================================================
\section{Palabras con varios significados}
%=====================================================================

Estas son las palabras que más confusiones producen en una defensa. Conviene
decir siempre cuál de los sentidos se usa.

{\small
\begin{longtable}{@{}p{0.17\textwidth}p{0.79\textwidth}@{}}
\toprule
\textbf{Palabra} & \textbf{Sentidos distintos}\\
\midrule
\endhead
completo / completitud & (1) \emph{funcional}: una base expresa toda tabla. (2) \emph{del cálculo}: lo verdadero es demostrable (Cap.\ 11). (3) \emph{de Gödel} (1929): la misma idea para lógica de primer orden, fuera de la tesis. (4) La \emph{incompletitud} de Gödel (1931) es otra cosa: habla de la aritmética, no de la lógica proposicional.\\
corrección & (1) De un \emph{cálculo}: lo demostrable es verdadero (solidez, Ap.\ A). (2) De un \emph{algoritmo}: cumple su especificación (Cap.\ 12).\\
base & (1) \emph{De conectivos}: los primitivos, \(\{\neg,\to\}\). (2) \emph{Caso base} de una inducción. (3) Base de la confianza: el núcleo y los axiomas.\\
normal & (1) \emph{Forma normal \(\beta\)}: sin redexes. (2) \emph{Forma normal disyuntiva}: disyunción de conjunciones. (3) \emph{Normalización de pruebas}: eliminar desvíos.\\
\(\vdash\) & (1) Derivabilidad: \(\Gamma\vdash\varphi\). (2) Tipado: \(\Gamma\vdash M:\sigma\). (3) En el asistente, la meta pendiente. La tesis lo advierte en §10.3.3.\\
reducción & (1) \emph{\(\beta\)-reducción}: un paso de cómputo. (2) \emph{Reducción al absurdo}: \(\neg\)I (intuicionista) o RAA (clásica). (3) \(\iota\)-reducción.\\
eliminación & (1) \emph{Regla de eliminación}: usar un conectivo (\(\to\)E). (2) \emph{Eliminador}: el recursor de un tipo inductivo.\\
abstracción & (1) \(\lambda x.M\). (2) Regla de tipado que la introduce.\\
tipo & (1) Clasificación de un término. (2) Bajo PAT, una proposición. (3) \emph{Tipo inductivo}: tipo definido por constructores.\\
consistente & (1) Un cálculo: no deriva \(\bot\). (2) Un conjunto de fórmulas: no deriva \(\bot\) a partir de él (Lindenbaum).\\
decidir & (1) Decidibilidad: hay un algoritmo. (2) \(\mathsf{decidir}(x\leq p)\): calcular el valor booleano de una comparación (Cap.\ 12).\\
\(\Rightarrow\) & (1) Implicación del metalenguaje. (2) Reducción paralela en la prueba de Church--Rosser (Def.\ 7.39).\\
\(\bot\) & (1) Fórmula falsa del cálculo. (2) A veces se usa para «Falso»; la tesis los distingue (Nota 11.1).\\
mínima & (1) \emph{Lógica mínima}: intuicionista sin ex falso. (2) \emph{Base mínima} de conectivos \(\{\neg,\to\}\): sigue siendo lógica clásica completa. La tesis lo aclara en el Cap.\ 9.\\
clásico / clásica & (1) \emph{Lógica clásica}: con LEM o RAA. (2) \emph{Resultado clásico}: conocido, de la literatura. (3) En Lean, \texttt{Classical}: el axioma de elección y sus consecuencias, como \texttt{Classical.em}.\\
inducción & Sobre \(\mathbb N\), estructural, sobre derivaciones, bien fundada: misma idea, distinto conjunto.\\
\bottomrule
\end{longtable}
}

%=====================================================================
\section{Tabla de símbolos}
%=====================================================================

{\small
\begin{longtable}{@{}p{0.2\textwidth}p{0.76\textwidth}@{}}
\toprule
\textbf{Símbolo} & \textbf{Lectura}\\
\midrule
\endhead
\(\Gamma\vdash\varphi\) & «\(\varphi\) se deriva de \(\Gamma\)»; \(\Provnd\) en el Cap.\ 11\\
\(\Gamma\models\varphi\) & «\(\varphi\) es consecuencia semántica de \(\Gamma\)»\\
\(\Gamma\vdash M:\sigma\) & «en \(\Gamma\), \(M\) tiene tipo \(\sigma\)»\\
\(\sem{\varphi}_v\) & valor de verdad de \(\varphi\) bajo la valuación \(v\)\\
\(v\models\Gamma\) & \(v\) hace verdaderas todas las fórmulas de \(\Gamma\)\\
\(\mathrm{Taut}(\varphi)\) & \(\varphi\) es tautología\\
\(\varphi\dashv\vdash\psi\) & demostrablemente equivalentes\\
\(\equiv\) & identidad sintáctica (en la tesis, módulo \(\alpha\))\\
\(=_\alpha,\ =_\beta\) & \(\alpha\)-equivalencia; \(\beta\)-conversión\\
\(\bred,\ \bmred\) & un paso de \(\beta\); cero o más pasos\\
\(M[x:=N]\) & sustitución de \(x\) por \(N\) en \(M\)\\
\(\mathrm{FV}(M)\) & variables libres de \(M\)\\
\(\lambda x.M\) & función de \(x\) con cuerpo \(M\)\\
\(\Lambda\alpha{:}*.M\) & abstracción sobre un tipo (polimorfismo)\\
\(\Pi x{:}A.B\) & producto dependiente (\(\forall\))\\
\(*,\ \square\) & sorts: tipos y kinds\\
\(\sigma\to\tau\) & tipo de funciones; implicación\\
\(\neg,\to,\wedge,\vee,\bot\) & no, implica, y, o, falso\\
\(\Gamma,\varphi\) & \(\Gamma\cup\{\varphi\}\)\\
\(a::r\) & lista con cabeza \(a\) y cola \(r\)\\
\(\ell_1\concat\ell_2\) & concatenación\\
\(|\ell|\) & longitud de la lista\\
\(\ell_1\sim\ell_2\) & \(\ell_1\) es permutación de \(\ell_2\)\\
\(\Gamma_v(A)\) & contexto de literales de la fila \(v\)\\
\(\Theta(\cdot)\) & orden de crecimiento (complejidad)\\
\(::=\) & «se define por la gramática»\\
\(:=\) & «se define como»\\
\bottomrule
\end{longtable}
}

\end{document}
